% 04_implementation.tex — converted from content/book.md, "4. Implementation
% Details" (Stage 9.3). Code snippets use plain verbatim (no extra packages);
% citations mirror the [6]/[7] markers in the cleaned Markdown.

\chapter{Implementation Details}

\section{Agent Definitions and Personalities}

Each agent is created with a role, goal, and backstory. The example below
shows how an agent is configured in CrewAI. It uses the same model our real
run used, \texttt{gemini/gemini-2.5-flash}, rather than naming any specific
commercial product as a project result:

\begin{verbatim}
from crewai import Agent, LLM

llm = LLM(model="gemini/gemini-2.5-flash", temperature=0.2)

planner_agent = Agent(
    role="Content Planner",
    goal="Develop a clear, hierarchical outline for a technical LaTeX book.",
    backstory=(
        "You are an experienced technical editor. You break broad topics "
        "into logical chapters and sections and keep the table of contents "
        "coherent."
    ),
    llm=llm,
    verbose=True,
    allow_delegation=False,
)
\end{verbatim}

The writer, reviewer, and reference-curator agents follow the same pattern
with their own goals and backstories. The model is supplied from
configuration, and the API key is read from the environment only --- it is
never written into the source.

\section{Custom Tool Integration}

Tools extend agents beyond plain text generation. A realistic example for
this project is a \LaTeX{} validation tool that checks whether generated
\LaTeX{} compiles. The following is illustrative pseudocode for such a tool:

\begin{verbatim}
# latex_validation_tool.py (illustrative)
import subprocess
import tempfile
from pathlib import Path

def validate_latex(latex_code: str) -> str:
    """Attempt to compile a LaTeX snippet and report errors."""
    with tempfile.TemporaryDirectory() as tmp:
        tex = Path(tmp) / "snippet.tex"
        tex.write_text(
            "\\documentclass{article}\\begin{document}\n"
            f"{latex_code}\n\\end{{document}}\n",
            encoding="utf-8",
        )
        result = subprocess.run(
            ["pdflatex", "-interaction=nonstopmode",
             "-output-directory", tmp, str(tex)],
            capture_output=True,
            text=True,
        )
        if result.returncode == 0:
            return "Validation successful: no compilation errors."
        return "Validation failed: see compiler output for details."
\end{verbatim}

This shows the general idea of giving an agent a deterministic check it can
call. In the current project the deterministic checks live in plain Python
modules (for example the content-QA scanner) rather than being driven by the
LLM, which keeps the validation layer fully offline and reproducible.

\section{Task Design and Execution Flow}

Tasks are the atomic units of work. A simplified sequence for this project:

\begin{enumerate}
  \item \textbf{Plan outline} --- agent: planner; input: the topic; output:
        a structured outline.
  \item \textbf{Draft content} --- agent: writer; input: the outline;
        output: the draft.
  \item \textbf{Review} --- agent: reviewer; input: the draft; output:
        structured feedback.
  \item \textbf{Curate references} --- agent: reference curator; input:
        draft and review; output: the reference list.
\end{enumerate}

Table~\ref{tab:agents} summarises the agents, their roles, and example
tools.

\begin{table}[htbp]
  \centering
  \caption{Agents, roles, and example tools.}
  \label{tab:agents}
  \begin{tabular}{llll}
    \toprule
    Agent & Primary role & Example tasks & Example tools \\
    \midrule
    Planner & Structure the book & Outline generation, section split &
      OutlineTool \\
    Writer & Generate the prose & Drafting chapters and sections &
      WritingTool \\
    Reviewer & Quality assurance & Coverage and consistency checks &
      ReviewTool \\
    Reference curator & Collect citations & Building the reference list &
      ReferenceTool \\
    \bottomrule
  \end{tabular}
\end{table}

\section{Handling Multilingual Content}

The pipeline is designed to support content in more than one language, which
matters for documentation that mixes scripts. The relevant \LaTeX{} support
comes from packages such as \texttt{babel} or \texttt{polyglossia}, and for
bidirectional text the \texttt{bidi} package~\cite{ctan_bidi}, usually
compiled with XeLaTeX~\cite{tug_xetex} or LuaLaTeX for good Unicode and
right-to-left handling.

\paragraph{Hebrew--English bidirectional demonstration.} The following
passage demonstrates bidirectional typesetting: the text direction switches
automatically inside the right-to-left environment. This paragraph mixes
English and Hebrew text.

\begin{hebrew}
זוהי דוגמה למקטע הממחיש יצירת תוכן דו-לשוני.
\end{hebrew}

The English text flows left-to-right while the Hebrew sentence (meaning
``this is an example of a section demonstrating the creation of bilingual
content'') flows right-to-left, both within the same document ---
demonstrating correct bidirectional integration.
